%=====================================================================
%  SHARED PREAMBLE — Quantitative Reasoning LAB 413 Beamer theme
%  Paste this block VERBATIM at the top of every LabN_beamer.tex
%  (keep each deck self-contained; do NOT \input this file).
%  Compile: pdflatex LabN_beamer.tex   (run twice for footline counters)
%=====================================================================
\documentclass[aspectratio=169,11pt]{beamer}

\usepackage[utf8]{inputenc}
\usepackage[T1]{fontenc}
\usepackage{lmodern}
\usepackage{amsmath,amssymb}
\usepackage{graphicx}
\usepackage{booktabs}
\usepackage{tcolorbox}
\tcbuselibrary{skins}

\usetheme{default}
\usecolortheme{default}
\useinnertheme{circles}

\definecolor{qrnavy}{RGB}{20,44,90}
\definecolor{qrblue}{RGB}{38,80,150}
\definecolor{qrgray}{RGB}{90,95,105}
\definecolor{qrlight}{RGB}{234,239,247}
\definecolor{qrred}{RGB}{190,60,60}

\setbeamercolor{structure}{fg=qrnavy}
\setbeamercolor{frametitle}{fg=white,bg=qrnavy}
\setbeamercolor{title}{fg=qrnavy}
\setbeamercolor{block title}{fg=white,bg=qrblue}
\setbeamercolor{block body}{bg=qrlight}
\setbeamercolor{itemize item}{fg=qrblue}
\setbeamercolor{itemize subitem}{fg=qrblue}
\setbeamerfont{frametitle}{series=\bfseries}
\setbeamertemplate{navigation symbols}{}

\setbeamertemplate{frametitle}{%
  \nointerlineskip
  \begin{beamercolorbox}[wd=\paperwidth,ht=3.2ex,dp=1.2ex,leftskip=0.3cm]{frametitle}%
    \strut\insertframetitle\strut
  \end{beamercolorbox}%
}

% Footline: course | lab N | slide number  (edit "Lab N" per deck)
\newcommand{\labnum}{Lab 10}   % <-- SET THIS to the lab number, e.g. \renewcommand{\labnum}{Lab 3}
\setbeamertemplate{footline}{%
  \leavevmode%
  \hbox{%
    \begin{beamercolorbox}[wd=.5\paperwidth,ht=2.4ex,dp=1ex,leftskip=.3cm]{author in head/foot}%
      \usebeamerfont{author in head/foot}\color{qrgray}\small Quantitative Reasoning -- LAB 413%
    \end{beamercolorbox}%
    \begin{beamercolorbox}[wd=.5\paperwidth,ht=2.4ex,dp=1ex,rightskip=.3cm]{title in head/foot}%
      \usebeamerfont{title in head/foot}\color{qrgray}\small\hfill \labnum \quad\insertframenumber\,/\,\inserttotalframenumber%
    \end{beamercolorbox}}%
}

% Number badge (01, 02, ...) used on concept slides:  \numbadge{01}
\newtcbox{\numbadge}{on line,boxsep=2pt,left=4pt,right=4pt,top=2pt,bottom=2pt,
  colback=qrnavy,colframe=qrnavy,fontupper=\bfseries\color{white},arc=2pt}

% Key-concepts callout:  \begin{keybox} ... \end{keybox}   or  \begin{keybox}[Scenario] ...
\newtcolorbox{keybox}[1][Key Concepts]{colback=qrlight,colframe=qrnavy,
  fonttitle=\bfseries,title={#1},coltitle=white,arc=2pt,left=4pt,right=4pt,
  top=3pt,bottom=3pt,boxrule=0.6pt}

%---------------------------------------------------------------------
%  Title data
%---------------------------------------------------------------------
\title{\bfseries Inference for Regression}
\subtitle{Week 11 Recap \textbullet\ Lab 10}
\author{Instructor: Subin Na}
\institute{Quantitative Reasoning -- LAB 413}
\date{November 12, 2025}

\begin{document}

%---------------------------------------------------------------------
%  Title slide
%---------------------------------------------------------------------
{\setbeamertemplate{footline}{}
\begin{frame}
  \vfill
  {\color{qrgray}\small Quantitative Reasoning -- LAB 413}\\[0.6em]
  {\usebeamerfont{title}\color{qrnavy}\huge\bfseries Inference for Regression,\\[0.15em]Multivariate Regression \& Interaction Terms\par}
  \vspace{0.4em}
  {\color{qrblue}\large Week 11 Recap \ \textbullet\ \ Lab 10}\par
  \vspace{1.6em}
  {\color{qrgray} Instructor: \textbf{Subin Na} \hfill November 12, 2025}
  \vfill
\end{frame}
}

%---------------------------------------------------------------------
%  Today's Plan
%---------------------------------------------------------------------
\begin{frame}{Today's Plan}
  \begin{columns}[T]
    \begin{column}{0.5\textwidth}
      \textbf{\color{qrnavy}1. Quick Recap}
      \begin{itemize}
        \item Inference for Proportions (review)
        \item Inference for Regression
        \item OLS Assumptions
        \item Multiple Regression Concept
      \end{itemize}
    \end{column}
    \begin{column}{0.5\textwidth}
      \vspace{2.2em}
      \textbf{\color{qrnavy}2. RStudio Practice}
    \end{column}
  \end{columns}
\end{frame}

%---------------------------------------------------------------------
%  01 Inference for Proportions (1)
%---------------------------------------------------------------------
\begin{frame}{\numbadge{01}\ \ Inference for Proportions (1)}
  \begin{itemize}
    \item We use sample proportions $\hat{p}$ to estimate the population
          proportion $p$.
    \item Because samples vary, $\hat{p}$ has its own \textbf{sampling
          distribution} $\rightarrow$ roughly Normal if:
      \begin{itemize}
        \item Random sample
        \item Expected successes $\geq 10$, failures $\geq 10$
      \end{itemize}
  \end{itemize}
  \vspace{0.4em}
  \begin{keybox}[Confidence Interval (CI)]
    A \textbf{range of plausible population proportions}:
    \begin{equation*}
      \hat{p} \;\pm\; z^{*}\,\sqrt{\dfrac{\hat{p}(1-\hat{p})}{n}}
    \end{equation*}
  \end{keybox}
\end{frame}

%---------------------------------------------------------------------
%  01 Inference for Proportions (2) - worked example setup
%---------------------------------------------------------------------
\begin{frame}{\numbadge{01}\ \ Inference for Proportions (2)}
  \begin{keybox}[Example]
    In a sports study, \textbf{73 out of 150} athletes reported sleeping more
    than 8 hours the previous night. Compute and interpret a \textbf{90\%
    confidence interval} for this proportion.
  \end{keybox}
  \vspace{0.8em}
  \begin{itemize}
    \item Point estimate: $\hat{p} = \dfrac{73}{150} \approx 0.487$.
    \item Critical value for a 90\% CI: $z^{*} = 1.645$.
  \end{itemize}
\end{frame}

%---------------------------------------------------------------------
%  Worked example - solution
%---------------------------------------------------------------------
\begin{frame}{Inference for Proportions: Solution}
  Plug in and solve:
  \vspace{0.4em}
  \begin{align*}
    \hat{p} \;\pm\; z^{*}\,\sqrt{\dfrac{\hat{p}(1-\hat{p})}{n}}
      &= 0.487 \;\pm\; 1.645\,\sqrt{\dfrac{(0.487)(1-0.487)}{150}}\\[0.5em]
      &= 0.487 \;\pm\; 0.067
  \end{align*}
  \vspace{0.2em}
  \begin{center}
    \begin{tcolorbox}[colback=qrlight,colframe=qrnavy,arc=2pt,boxrule=0.6pt,width=0.7\textwidth,halign=center]
      90\% CI is $(0.420,\ 0.554)$
    \end{tcolorbox}
  \end{center}
\end{frame}

%---------------------------------------------------------------------
%  02 Inference for Regression
%---------------------------------------------------------------------
\begin{frame}{\numbadge{02}\ \ Inference for Regression}
  \begin{itemize}
    \item Regression estimates $(\hat{\beta}_0,\hat{\beta}_1)$ are statistics
          computed from samples.
    \item Each has a \textbf{sampling distribution} $\rightarrow$ uncertainty
          measured by the standard error (SE).
  \end{itemize}
  \vspace{0.3em}
  \begin{keybox}[Key Questions]
    \begin{itemize}
      \item Is the relationship \textbf{real} or due to chance?
      \item How \textbf{precise} is our slope estimate?
    \end{itemize}
  \end{keybox}
  \vspace{0.3em}
  \begin{itemize}
    \item $95\%$ CI $\approx \hat{\beta} \pm 1.96 \times \text{SE}$.
    \item Reject $H_0$ if $p < 0.05$ ($\beta$ significantly different from $0$).
  \end{itemize}
\end{frame}

%---------------------------------------------------------------------
%  03 OLS Assumptions
%---------------------------------------------------------------------
\begin{frame}{\numbadge{03}\ \ OLS Assumptions}
  \begin{itemize}
    \item \textbf{Linearity} --- relationship between $X$ and $Y$ is linear.
    \item \textbf{Random sampling} --- observations are independent and random.
    \item \textbf{Exogeneity} --- errors uncorrelated with predictors.
    \item \textbf{Homoskedasticity} --- constant error variance.
    \item \textbf{Variation in $X$} --- predictors are not constant.
  \end{itemize}
  \vspace{0.6em}
  \begin{center}
    {\color{qrgray}\small (Violating these assumptions $\Rightarrow$ biased or
    inefficient estimates.)}
  \end{center}
\end{frame}

%---------------------------------------------------------------------
%  04 Multiple Regression Concept
%---------------------------------------------------------------------
\begin{frame}{\numbadge{04}\ \ Multiple Regression Concept}
  \begin{keybox}[Model]
    \begin{equation*}
      Y = \beta_0 + \beta_1 X_1 + \beta_2 X_2 + \dots + \varepsilon
    \end{equation*}
  \end{keybox}
  \vspace{0.4em}
  \begin{itemize}
    \item Adds more information and \textbf{controls for confounding}.
    \item Each coefficient is interpreted \emph{holding other variables constant}.
    \item Helps separate true effects from spurious associations.
    \item Addresses \textbf{omitted variable bias} if relevant controls are included.
  \end{itemize}
\end{frame}

%---------------------------------------------------------------------
%  Transition to RStudio
%---------------------------------------------------------------------
{\setbeamertemplate{footline}{}
\begin{frame}
  \vfill
  \centering
  {\color{qrnavy}\Huge\bfseries RStudio Practice}\\[0.8em]
  {\color{qrblue}\large Proportion Tests \& Inference for Regression}\\[0.4em]
  {\color{qrgray} See \texttt{Lab10.Rmd} / \texttt{Lab10\_1112.R}}
  \vfill
\end{frame}
}

\end{document}
